\documentclass[handout]{beamer}
\usepackage{array}
\usepackage{german}
\usepackage{graphicx}
\usepackage[utf8]{inputenc}
\usepackage[T1]{fontenc}
\mode<beamer>{%
\usetheme{Copenhagen}
}
\usepackage[orientation=landscape,size=a3,debug,scale=1.9]{beamerposter}
\title[]{}
\begin{document}
\begin{frame}
\frametitle{%\hspace{0pt plus 1 filll}
MathSem 2026: Analysis und Algebra}
\vspace*{-0.1cm}
\begin{columns}[t,onlytextwidth]
\begin{column}{0.37\textwidth}
\begin{center}
\includegraphics[width=\hsize]{../cover/buchcover.png}
\end{center}
%\vskip 0.2cm
\bigskip
Umschlagbild: Visualizierung einer meromorphen Funktion
\bigskip
\bigskip

$\text{xvi} + 754$ Seiten, gebunden.
Erscheint~im~Herbst~2026.\\
Anfragen~an~Prof.~Dr.~Andreas Müller,\\
{\texttt{andreas.mueller@ost.ch}}
\bigskip
\bigskip
\bigskip
%\vskip 1.2cm
\end{column}
%
\begin{column}{0.60\textwidth}
\begin{description}
\item[Teil 1:] Grundlagen
\begin{enumerate}
\item Einleitung
\item Komplexe Zahlen und Funktionen
\item Holomorphe Funktionen
\item Komplexe Integration
\item Analytische Funktionen
\item Differentialgleichungen
\item Integration rationaler Funktionen
\item Integration in geschlossener Form
\end{enumerate}
\item[Teil 2:] Anwendungen und weiterführende Themen
\begin{enumerate}
\setcounter{enumi}{8}
\item
\textit{Timon Gnehm und Pascal Hirzel}: Komplexe Zahlen und geometrische Konstruktionen
\item
\textit{Tom Dawson}: Möbius-Transformationen
\item
\textit{Shaarujan Kamalanathan}: Das seltsame Problem der geradlinigen Bewegung
\item
\textit{Nicola Dall'Acqua}: Julia-Mengen und die Mandelbrot-Menge
\item
\textit{Dominik Castelberg}: Der jordansche Kurvensatz
\item
\textit{Kai Erdin}: Quaternionen
\item
\textit{Malenka Hossli} und \textit{Andrea Studer}: Zweidimensionale Elektrostatik
\item
\textit{Jack Dumovich} und \textit{Roman Meyer}: Die Auftriebsformel von Kutta-Joukowski
\item
\textit{Stephanie Märklin}: Aircraft Flight Dynamics und komplexe Zahlen
\item
\textit{Andri Sprecher} und \textit{Fabian Suter}: Fresnel-Integrale
\item
\textit{Michael Schmid}: Padé-Approximation
\item
\textit{Selina Malacarne}: Complex Step Derivative
\item
\textit{Damien Flury}: Hauptwert eines divergenten Integrals
\item
\textit{Jan Klarer} und \textit{Manuel Kuhn}: Das dritte keplersche Gesetz
\item
\textit{Noel Karlsson} und \textit{Simon Widmer}: Die Laplace-Inversion: Vom Bromwich-Integral zur Talbot-Kontour
\item
\textit{Lukas Krüdewagen}: Der Weyl-Algorithmus
\item
\textit{Simon Köpfli} und \textit{Florian von Wyl}: Produktentwicklungen
\item
\textit{Jan Gachnang} und \textit{Etienne Schafflützel}: Das Basel-Problem
\item
\textit{Lukas Buchli} und \textit{Roman Weber}: Die Gamma-Funktion
\item
\textit{Gian Cavegn}: Bessel-Funktionen
\item
\textit{Jannis Gull}: Zeta-Funktion und Hankel-Contour
\item
\textit{Melissa Haumüller}: Die riemannsche Zeta-Funktion und die Summer der natürlichen Zahlen
\item
\textit{Nathanael Fässler}: Der Buchberger-Algorithmus
\item
\textit{Jakob Gierer}: Gemeinsame Nullstellen und Resultante
\item
\textit{Baris Catan}: Elliptische Integrale
\end{enumerate}
\end{description}
\end{column}
\begin{column}{0.01\textwidth}
\llap{\raisebox{-3cm}{\includegraphics{qrcode.pdf}}}
\end{column}
\end{columns}
\end{frame}
\end{document}
